\section{Kesimpulan}

Pengujian perangkat lunak menunjukkan bahwa instrumen menghasilkan penandaan gerak paralel dan resolusi nada penuntun pada kasus uji yang digunakan. Sampel paralel artifisial menghasilkan skor 0,0000, sedangkan BWV 66.6 menghasilkan skor 0,7568 pada konfigurasi tonika C. Hasil tersebut menggambarkan perilaku instrumen pada kedua sampel, dengan batas interpretasi yang dijelaskan pada Bab IV.

Skor menggunakan jumlah penandaan yang dinormalisasi terhadap peristiwa sopran. Karena satu waktu dapat memuat beberapa penandaan, indeks ini tidak identik dengan proporsi momen harmonis tanpa pelanggaran. Interpretasi musikal bergantung pada definisi interval, representasi waktu, pemetaan suara, dan konteks tonal.

\section{Saran}

Pengembangan instrumen perlu memperhatikan representasi nada yang ditahan, ketepatan kuantisasi, serta pemisahan interval kuart dan kuint. Pemeriksaan resolusi nada penuntun dapat diperluas ke suara alto, tenor, dan bas dengan mempertimbangkan fungsi harmonisnya. Anotasi contoh secara independen dapat digunakan untuk mengukur kesalahan deteksi.

Perbandingan model dengan tugas dan data latih yang lebih terkendali dapat membantu memisahkan pengaruh arsitektur dari pengaruh representasi, pelatihan, dan prosedur generasi.
